% Generated draft: requires mathematical and source review.
\documentclass[11pt]{article}
\usepackage[margin=25mm]{geometry}
\usepackage{fontspec}
\setmainfont{Latin Modern Roman}
\usepackage{amsmath,amssymb,mathtools,unicode-math}
\setmathfont{Latin Modern Math}
\usepackage{enumitem,needspace,xurl,hyperref}
\setlength{\parindent}{0pt}
\setlength{\parskip}{5pt}
\setlength{\emergencystretch}{3em}
\newcommand{\N}{\mathbb{N}}
\newcommand{\Z}{\mathbb{Z}}
\newcommand{\Q}{\mathbb{Q}}
\newcommand{\R}{\mathbb{R}}
\newcommand{\C}{\mathbb{C}}
\newcommand{\F}{\mathbb{F}}
\newcommand{\A}{\mathbb{A}}
\newcommand{\PP}{\mathbb{P}}
\newcommand{\E}{\mathbb{E}}
\newcommand{\eps}{\varepsilon}
\newcommand{\dd}{\,\mathrm{d}}
\begin{document}
\section*{UnsolvedMath ID: 30001826}
\textbf{Source problem code:} catalogue entry
\subsection{Definitions and mathematical statement}
Write \(M=X\setminus D\) for the complex surface appearing in the source record. The intended objects are compactifiable gravitational instantons, but the record does not specify the hypotheses on \(X\) and \(D\), the permitted compactifications, or a definition of a gravitational-instanton metric. These choices must remain unresolved; in particular, completeness, curvature integrability, asymptotic conditions, and compatibility with the given complex structure cannot be silently supplied.

For clarity, a hyperkähler structure on a smooth real four-manifold consists of a Riemannian metric \(g\) and parallel complex structures \(I,J,K\), each orthogonal for \(g\), satisfying
\[
I^2=J^2=K^2=IJK=-\operatorname{id}.
\]
For such a structure, a hyperkähler rotation replaces the complex structure by
\[
I_{a,b,c}=aI+bJ+cK,
\qquad (a,b,c)\in\R^3,\qquad a^2+b^2+c^2=1.
\]
A biholomorphism is a holomorphic bijection with holomorphic inverse.

\textbf{Classification problem.} Classify the complex surfaces \(X\setminus D\) carrying gravitational-instanton metrics, up to biholomorphism and hyperkähler rotation. The original formulation asks whether every such surface becomes biholomorphic, after hyperkähler rotation, to one of the known examples.

Schematically, if \(\mathcal C\) denotes the intended class of admissible complements carrying such metrics and \(\mathcal K\) denotes the intended collection of known examples, the proposed exhaustiveness assertion is
\[
\forall M\in\mathcal C\;\exists N\in\mathcal K:\quad
M\text{ is biholomorphic to }N\text{ after hyperkähler rotation}.
\]
This is a schematic statement requiring further source clarification: neither \(\mathcal C\) nor \(\mathcal K\) is specified sufficiently in the record to make it a fully determined conjecture. The record also does not settle whether the rotation assertion concerns every compatible gravitational-instanton metric on \(M\), or the existence of a suitable such metric and hyperkähler structure.
\subsection{Short English statement}
Classify compactifiable complex surfaces carrying gravitational-instanton metrics, allowing biholomorphisms and changes of complex structure by hyperkähler rotation. The source asks whether known examples exhaust this class, but does not specify the admissible compactifications, metric hypotheses, or comparison examples. The supplied open label and literature assessment require independent review.
\subsection{Sources}
\begin{itemize}
\item Local UnsolvedMath registry: permanent ID \(30001826\), title \emph{Compactifiable Gravitational Instantons}, source code \texttt{catalogue entry}, in the supplied Oberwolfach Reports Open Problems collection. Project registry locations are \texttt{data/unsolvedmath/problems.json} and \texttt{catalog/problems.json}; this draft uses the supplied complete record. No individual proposer is identified; the supplied proposed year is 2011.
\item Supplied report citation: \emph{Differentialgeometrie im Großen (2012)}, Oberwolfach Reports, DOI 10.4171/owr/2011/33, pp.\ 1857--1912. Original report link, also supplied as a literature source: \url{https://doi.org/10.4171/owr/2011/33}. The report is described as documenting workshop discussions and open problems. The supplied title year 2012 differs from the proposed year 2011 and the year component of the DOI; this discrepancy is unresolved.
\item Additional supplied literature source, labeled ``recent gravitational-instanton classification'': \url{https://arxiv.org/abs/2310.13197}. No authors, precise title, theorem, or explanation of its scope are supplied.
\item Supplied literature assessment dated exactly \texttt{2026-08-21}: the record labels the problem open and reports that the extracted classification conjecture for all compactifiable gravitational instantons is broader than the cited examples, with no complete proof or counterexample in the checked literature. This is supplied dated curation derived from the OWR record, not an independently verified literature assessment. The supplied background explicitly calls for independent verification of time-sensitive status and cited claims.
\item The supplied statement status is \texttt{corrected\_verified}, with a supplied claim of manual semantic, typographic, and LaTeX review. These are registry annotations, not independent verification of the mathematical formulation, the original report, or the open status.
\end{itemize}
\end{document}
